\section{One-Fluid 化の数学的メカニズム：1次元モデルによる検証}
\label{app:onefluid_1d}
% ═══════════════════════════════════════════════════════════════════════════════

本稿 §\ref{sec:two_to_one} で述べた Two-Fluid から One-Fluid への変換がなぜ成り立つのかを，
1次元の簡単な拡散方程式で示す．

\subsection*{問題設定}

1次元領域 $\Omega = (-1, 1)$ で，$x=0$ に界面があるとする．
左域 $\Omega^- = (-1,0)$（液相，係数 $\mu_l$）と
右域 $\Omega^+ = (0,1)$（気相，係数 $\mu_g$）で
それぞれ拡散方程式が成立し，界面条件を持つ：
\begin{align}
  \frac{d}{dx}\!\left(\mu_l \frac{du}{dx}\right) &= 0 \quad (x \in \Omega^-) \\
  \frac{d}{dx}\!\left(\mu_g \frac{du}{dx}\right) &= 0 \quad (x \in \Omega^+)
\end{align}
\textbf{界面条件（$x=0$）：}
\begin{align}
  [u]_0 &= u(0^+) - u(0^-) = 0 \quad\text{（速度の連続性）} \\
  \left[\mu\frac{du}{dx}\right]_0
  &= \mu_g u'(0^+) - \mu_l u'(0^-)
   = f_\sigma \quad\text{（応力のジャンプ = 界面力）}
\end{align}

\subsection*{One-Fluid 化}

$\mu(x) = \mu_l H(-x) + \mu_g H(x)$（ヘヴィサイド関数で補間した係数）として，
全領域で1本の方程式を書く：
\begin{equation}
  \frac{d}{dx}\!\left(\mu(x) \frac{du}{dx}\right) = f_\sigma\,\delta(x)
  \label{eq:app_1d_onefluid}
\end{equation}
ここで右辺の $\delta(x)$ は界面 $x=0$ に集中するディラックデルタ関数である．

\subsection*{数学的正当性：積分による検証}

式 \eqref{eq:app_1d_onefluid} を区間 $(-\varepsilon, +\varepsilon)$（$\varepsilon \to 0$）で積分する：
\[
  \int_{-\varepsilon}^{+\varepsilon} \frac{d}{dx}\!\left(\mu\frac{du}{dx}\right) dx
  = \left[\mu\frac{du}{dx}\right]_{-\varepsilon}^{+\varepsilon}
  = \mu_g u'(0^+) - \mu_l u'(0^-)
\]
右辺：
\[
  \int_{-\varepsilon}^{+\varepsilon} f_\sigma\,\delta(x)\,dx = f_\sigma
\]
両辺を比較すると，応力ジャンプ条件が自動的に再現されることがわかる．

また，$\mu(x)$ が $x = 0$ で不連続であることを，弱形式（分布関数的な意味）で処理することにより，
$x=0$ での方程式 \eqref{eq:app_1d_onefluid} は
速度連続条件・応力ジャンプ条件と等価になる．

\subsection*{結論と3次元への一般化}

One-Fluid 方程式に右辺のデルタ関数項を加えることで，
界面条件が体積方程式の中に「吸収」される．
3次元 Navier-Stokes 方程式における CSF モデル $\bm{f}_\sigma = \sigma\kappa\nabla\psi$ は，
この1次元モデルの表面張力（法線応力のジャンプ）を多次元・曲面界面に一般化したものである．\qed

% ═══════════════════════════════════════════════════════════════════════════════
